\chapter{Declaración de uso de Inteligencia Artificial}\label{cap:uso-ia}

De acuerdo con las directrices de la asignatura, se declara de forma transparente
el uso de herramientas de inteligencia artificial generativa durante el desarrollo
de este proyecto.

\section{Generación del código}
La mayor parte del código del proyecto (notebooks de Zeppelin, configuración
Docker y scripts auxiliares) ha sido generada con asistencia de \textbf{Claude}
(Anthropic), un modelo de lenguaje conversacional. Esta decisión se tomó por
la falta de experiencia previa del equipo con las tecnologías específicas del
stack (PySpark, Spark SQL, Folium, GeoJSON).

No obstante, \textbf{todo el código ha sido revisado, comprendido y validado
por los autores} antes de su incorporación al proyecto. Las decisiones de
diseño, la selección de consultas y la estructura del flujo analítico
responden a criterios propios del equipo.

La conversación utilizada para la generación del código está disponible en el
siguiente enlace:

\begin{center}
\url{https://claude.ai/share/231f84af-d261-4aad-b579-fa9fdaf06f2e}
\end{center}

\section{Redacción de la memoria}
La memoria LaTeX ha sido desarrollada con asistencia del IDE
\textbf{Antigravity} (Google DeepMind), que integra un asistente de IA para
edición de código y documentación. Dado que Antigravity es una herramienta
local de escritorio, \textbf{no existe posibilidad de compartir o exportar
el historial de conversación} de la misma forma que con un chatbot web.

La IA se ha utilizado en parte para:
\begin{itemize}
    \item estructurar y redactar las secciones de la memoria;
    \item desarrollar los análisis textuales de las gráficas y darles un
    significado más profundo en el contexto de movilidad urbana;
    \item mantener coherencia entre capítulos y asegurar que la documentación
    refleja fielmente el estado real del repositorio.
\end{itemize}

En todos los casos, el trabajo se ha realizado \textbf{bajo supervisión
directa de los autores}, que han dirigido, corregido y aprobado cada
modificación.

\section{Decisiones propias del equipo}
Es importante destacar que las siguientes decisiones son \textbf{íntegramente
de los autores}, sin intervención de la IA:

\begin{itemize}
    \item la \textbf{elección de las gráficas} incluidas en la memoria, que
    han sido seleccionadas por considerarlas las más interesantes y
    representativas del análisis;
    \item los \textbf{contenidos y temas analizados}, elegidos por su
    conveniencia y relevancia para responder a las preguntas de movilidad
    planteadas;
    \item la \textbf{arquitectura del proyecto} (Docker + Spark + Zeppelin),
    propuesta por el equipo docente y adaptada por los autores;
    \item la \textbf{interpretación final de los resultados} y las conclusiones
    extraídas del análisis.
\end{itemize}

\section{Compromiso de honestidad}
Los autores asumen la responsabilidad íntegra sobre el contenido del proyecto
y la memoria. El uso de herramientas de IA ha sido un medio para acelerar el
desarrollo, no un sustituto del criterio técnico y académico propio.
